% ==============================================================================
\section{Problem Statement}
\label{sec:problem}
% ==============================================================================

% ------------------------------------------------------------------------------
\subsection*{Operational overview}
% ------------------------------------------------------------------------------

A maintenance hangar consists of a finite number of parking positions
arranged in space.  Because of the physical geometry of the facility,
some positions are \emph{rear} positions whose only access path runs
across one or more \emph{front} positions.  An aircraft parked at a
front position therefore physically obstructs the entry and exit of
any aircraft at the rear positions behind it.  Whenever a rear-position
aircraft has to enter or leave the hangar while a front aircraft is
being serviced, the front aircraft must be temporarily towed out of
the way and returned to its slot once the manoeuvre is complete.  Each
such manoeuvre is operationally expensive: it occupies the ground crew,
interrupts service, and carries a non-negligible risk of damage.

A fleet of aircraft must be serviced in the hangar over a planning
horizon.  For each aircraft the airline specifies the work to be
performed, an earliest date on which work can start, and a target date
by which it should be finished; finishing later than the target incurs
a delay cost that the operator wants to minimise.  The work itself is
an ordered list of \emph{jobs} with known individual durations and a
strict chain of precedences: a job cannot start until its predecessor
has finished.  Each job is also tagged as either \emph{interruptible},
meaning the ground crew can pause it briefly to clear the way for a
neighbouring movement and resume it shortly afterwards, or
\emph{non-interruptible}, meaning it must run from start to finish
without external disturbance.

A schedule consists of two decisions per aircraft: which position to
assign it to, and when to start each of its jobs.  Two aircraft sharing
a position must be serviced one after the other.  An aircraft assigned
to a rear position can always enter or leave when the relevant front
position is unoccupied at that instant.  When the front position is
occupied, only two situations make the manoeuvre possible: (i) the
front aircraft is between two consecutive jobs and the gap is wide
enough to absorb the tow-out and tow-back; or (ii) the front aircraft
is mid-job and that job is interruptible, in which case the job is
paused, the manoeuvre is performed, and the job is resumed at the
cost of some additional working time.  Any other attempt to access the
rear position is operationally infeasible.  The schedule must satisfy
all of the above and minimise a weighted combination of makespan,
total delay, and number of manoeuvres.

% ------------------------------------------------------------------------------
\subsection*{Formal definition}
% ------------------------------------------------------------------------------

The hangar is modelled as a directed acyclic graph $G = (\calP, \calA)$
where $\calP = \{p_1, \dots, p_P\}$ is the set of parking positions and
$\calA \subseteq \calP \times \calP$ is the set of \emph{blocking
arcs}: an arc $(p, p') \in \calA$ means that position $p$ (the
\emph{front}) obstructs the access path to position $p'$ (the
\emph{rear}).

Let $\calR = \{r_1, \dots, r_R\}$ be the fleet of aircraft to be
serviced.  Each aircraft $r \in \calR$ carries an ordered chain of
jobs
\[
  \calJ_r = (j^r_1, j^r_2, \ldots, j^r_{N_r}),
\]
in which $j^r_{k+1}$ may only start after $j^r_k$ has finished.  Let
$\calJ = \bigcup_{r \in \calR} \calJ_r$ be the global job set.  Each
job $j \in \calJ$ has a nominal duration $D_j > 0$ and an
interruptibility flag
\[
  I_j \in \{0, 1\}, \qquad I_j = 1 \iff j \text{ is interruptible.}
\]
The aggregate processing time of aircraft $r$ is
$T_r = \sum_{j \in \calJ_r} D_j$.  Each aircraft also has an earliest
start $E_r \ge 0$ and a target finish $L_r > E_r$.

Four positive temporal parameters complete the instance:
\begin{itemize}[leftmargin=2em, itemsep=0pt]
  \item $\varepsilon$: the towing time of an aircraft into or out of a
        parking slot.  Used as the minimum separation between two
        consecutive aircraft serviced at the same position.
  \item $\mu$: the minimum inter-job pause that must be left between
        two consecutive jobs of an aircraft when a neighbouring access
        manoeuvre takes place during that pause.  In its absence the
        gap may be zero.
  \item $\delta$: the additional working time added to an interruptible
        job each time it is paused to allow a neighbouring access
        manoeuvre.
\end{itemize}
All three parameters are specified per instance.

% ------------------------------------------------------------------------------
%\subsection*{Decision variables and basic constraints}
% ------------------------------------------------------------------------------

A solution assigns each aircraft $r \in \calR$ to exactly one position
$\pi(r) \in \calP$, and fixes a start time $s_j \ge 0$ for every job
$j \in \calJ$.  The finish time of a job is
\begin{equation}
  f_j \;=\; s_j + D_j + \delta\, \kappa_j,
  \label{eq:job_finish}
\end{equation}
where $\kappa_j \in \NZ$ counts how many times job $j$ was paused
during its execution to enable a neighbouring access manoeuvre (see
the access conditions below); $\kappa_j > 0$ requires $I_j = 1$.
Job-level start and finish times induce aircraft-level start and finish
times $s_r = s_{j^r_1}$ and $f_r = f_{j^r_{N_r}}$, and the per-job
delay $v^D_r = \max(0,\, f_r - L_r)$.

The basic feasibility requirements are
\begin{align}
  & \textstyle\sum_{p \in \calP} \mathbf{1}[\pi(r) = p] = 1
    && \forall r \in \calR, \label{eq:assign} \\
  & s_{j^r_{k+1}} \;\ge\; f_{j^r_k}
    && \forall r \in \calR,\ k = 1, \ldots, N_r{-}1, \label{eq:chain} \\
  & s_r \;\ge\; E_r
    && \forall r \in \calR, \label{eq:earlystart} \\
  & [s_r, f_r] \cap [s_{r'}, f_{r'}] = \emptyset \text{ and gap } \ge \varepsilon
    && \forall r \ne r' \in \calR \text{ with } \pi(r) = \pi(r'). \label{eq:sameposseq}
\end{align}
Equation~\eqref{eq:assign} enforces that every aircraft is assigned to
exactly one position; \eqref{eq:chain} chains each aircraft's jobs in
the prescribed order; \eqref{eq:earlystart} respects the earliest
start; and \eqref{eq:sameposseq} forbids overlapping occupancies at
the same position and leaves at least $\varepsilon$ time units
between the finish of one aircraft and the start of the next.

% ------------------------------------------------------------------------------
\subsection*{Access conditions on blocked positions}
\label{subsec:access}
% ------------------------------------------------------------------------------

Consider any arc $(p, p') \in \calA$ and an aircraft $r'$ assigned to
the rear position $p' = \pi(r')$.  Its two access instants are the
entry $\tau = s_{r'}$ and the exit $\tau = f_{r'}$.  Each of these
instants must fall into exactly one of three modes, depending on the
state of the front position $p$ at time $\tau$.

\paragraph{Mode (A) — front position vacant at $\tau$.}
If no aircraft is being serviced at $p$ at time $\tau$ --- formally,
no aircraft $r$ with $\pi(r) = p$ satisfies $s_r \le \tau \le f_r$ ---
the access requires no manoeuvre and incurs no extra time and no
movement count.

\paragraph{Mode (B) — front aircraft mid-stay, $\tau$ in an inter-job gap.}
If there exists an aircraft $r$ with $\pi(r) = p$ such that $\tau$
falls between two consecutive jobs $j^r_k$ and $j^r_{k+1}$ of $r$,
i.e.
\begin{equation}
  f_{j^r_k} \;\le\; \tau \;\le\; s_{j^r_{k+1}},
  \label{eq:mode_b_window}
\end{equation}
then $r$ is towed out, $r'$ accesses $p'$, and $r$ is towed back.  The
gap must be wide enough to accommodate the manoeuvre:
\begin{equation}
  s_{j^r_{k+1}} \;-\; f_{j^r_k} \;\ge\; \mu
    \quad\text{whenever any access instant uses this gap}.
  \label{eq:mode_b_gap}
\end{equation}
The event contributes $2$ to the movement count (tow-out plus
tow-back) and does not extend any job.

\paragraph{Mode (C) — front aircraft mid-job, job interruptible.}
If there exists an aircraft $r$ with $\pi(r) = p$ such that $\tau$
falls strictly within the execution of some job $j^r_k$ of $r$, i.e.
\begin{equation}
  s_{j^r_k} \;<\; \tau \;<\; f_{j^r_k},
  \label{eq:mode_c_window}
\end{equation}
the access is feasible only if $I_{j^r_k} = 1$.  In that case $j^r_k$
is paused, the manoeuvre is performed, $j^r_k$ resumes, and the
counter $\kappa_{j^r_k}$ of equation~\eqref{eq:job_finish} is
incremented by $1$.  The event contributes $2$ to the movement count.

\paragraph{Infeasibility.}
If~\eqref{eq:mode_c_window} holds but $I_{j^r_k} = 0$, the access
instant $\tau$ is infeasible: the schedule must either retime $r'$ so
that $\tau$ falls into Mode~(A) or Mode~(B), or reassign $r'$ to a
different position.

\paragraph{Independence of entry and exit.}
The mode of the entry $s_{r'}$ and the mode of the exit $f_{r'}$ are
determined independently, against potentially different front aircraft
and different jobs.  Each contributes its own movement count and, in
Mode~(C), its own increment of $\kappa$.

% ------------------------------------------------------------------------------
\subsection*{Objective}
% ------------------------------------------------------------------------------

Let $m = \max_{r \in \calR} f_r$ be the makespan and let $n$ be the
total number of movements summed over all blocking arcs and aircraft
pairs.  The objective is the weighted combination
\begin{equation}
  \label{eq:obj}
  \min \;\; W^M \cdot m
        \;+\; W^D \sum_{r \in \calR} v^D_r
        \;+\; W^S \cdot n,
\end{equation}
where $W^M, W^D, W^S \ge 0$ are user-specified weights expressing the
relative priority of throughput, on-time delivery, and ground-crew
effort.

The problem is $\mathcal{NP}$-hard: even without blocking arcs it
contains the single-machine scheduling problem with arbitrary release
dates and weighted completion times~\citep{Lenstra1977}, and the
blocking access conditions of \S\ref{subsec:access} further enlarge
the combinatorial structure.  The mixed-integer programming
formulation that captures these conditions is developed in
\S\ref{sec:milp}.
